\section{Open Questions}
\label{sec:open}

The following five questions are things a follow-up study --- ours or someone
else's --- would need to answer to convert this PARTIAL replication into a
full one, or to extend the paper's clinical relevance.

\begin{enumerate}

\item \textbf{Open-source MC portability.}
Does an open-source Geant4-DNA reimplementation of the paper's geometry,
physics, and chemistry stages reproduce the $0.171 \pm 0.003$~DSB/decay
primitive for \textsuperscript{64}Cu at \SI{0.25}{nm}, or is the number
sensitive to TOPAS-nBio-specific chemistry parameters (\textbullet OH
lifetime, reaction radii, backbone volume definition)?\\
\emph{Next steps:} stand up Geant4-DNA v11.2+ (\texttt{G4EmDNAPhysics\_option4},
\texttt{G4EmDNAChemistry\_option3}), reconstruct the Zhu~et~al.\ 2020
fractal-chromatin geometry from published parameters, place
\textsuperscript{64}Cu at 0.25 and 1.15~nm on a random base pair, run
\texttt{G4RadioactiveDecay} for 400k histories per point, score SSBs with the
paper's 0.13 \textbullet OH probability + 17.5~eV direct threshold, and
compare DSB/decay under an identical 10-bp rule.

\item \textbf{Endpoint-model sensitivity.}
How sensitive is the \SI{47.1e-3}{Bq/cell} equi-lethal
\textsuperscript{64}Cu activity to the choice of biological-endpoint model
(Humm--Charlton 1989 vs.\ MEDRAS repair-kinetics vs.\ LEM-IV vs.\
PARTRAC-derived survival)?\\
\emph{Next steps:} wrap the Table 1 DSB primitives in (a) MEDRAS
fast/slow-repair, (b) LEM-IV with CHO parameters, (c) a PARTRAC
Poisson-cluster survival mapping.  Report the spread; a factor-of-2 spread
would materially change any clinical dose claim.

\item \textbf{Construct chemistry.}
Does DSB/decay depend materially on the specific
\textsuperscript{64}Cu-labeled targeting compound
(\textsuperscript{64}Cu-DOTA, -NOTA, -ATSM, -cyclam, antibody conjugates),
and how does that propagate to equi-lethal activity for realistic
clinical constructs?\\
\emph{Next steps:} compile published subcellular biodistribution for
\textsuperscript{64}Cu-ATSM, -DOTA-TATE, and antibody constructs; convolve
with a distance-to-DNA kernel to produce construct-specific effective
DSB/decay; compare to the paper's naked-\textsuperscript{64}Cu number.

\item \textbf{Auger sub-keV database sensitivity.}
How sensitive is the reported $0.171$~DSB/decay to the Auger sub-keV
spectrum, given that ICRP-107, ENSDF, and Livermore-Auger disagree by
factors of order 2 on the softest lines?  Sub-keV Auger electrons have
ranges $<\SI{10}{nm}$ and can dominate the DSB count.\\
\emph{Next steps:} rerun the Q1 Geant4-DNA replica under (i) default ICRP-107,
(ii) an explicit BrIccEmis per-shell spectrum.  Report the delta.  If
$>\SI{15}{\percent}$, tag the paper's number as spectrum-dependent and add a
systematic error bar to Table~2.

\item \textbf{ML surrogate for radionuclide screening.}
Can an ML fit over TOPAS-nBio outputs across
\{\textsuperscript{125}I, \textsuperscript{123}I, \textsuperscript{111}In,
\textsuperscript{99m}Tc, \textsuperscript{64}Cu\} recover a compact,
physically-interpretable RBE parameterisation that predicts DSB/decay for
a new Auger emitter without a fresh MC run?\\
\emph{Next steps:} build a small feature set (per-nuclide LET moments, mean
$E_e$, Auger yield $>1$~keV, source-axis distance, $T_{1/2}$); target =
DSB/decay from Table 1; fit linear + gradient-boosted models with
leave-one-nuclide-out CV.  If LOO MAE $<\SI{15}{\percent}$, publish as a
screening tool for \textsuperscript{67}Cu, \textsuperscript{161}Tb,
\textsuperscript{225}Ac daughters, etc.

\end{enumerate}
